% !TeX program = lualatex
% halo:begin
% {
%   "version": 1,
%   "publish": true,
%   "course": "higher-algebra-i",
%   "section": "1.5",
%   "title": "高等代数（一）1.5：行列式计算与拉普拉斯展开",
%   "categories": ["study-notes"],
%   "tags": ["mathematics"],
%   "excerpt": "行列式的降阶与递推计算、多行（列）的拉普拉斯展开、分块三角行列式与叉积。"
% }
% halo:end
% 课堂日期：2026-09-29；1.5为笔记自身的主题编号。
% 来源：笔记 2026年9月29日.pdf；IMG_1561、1563、1564、1565、1566、1568.HEIC。
% 与1.4衔接：不重复单行列展开的定义与证明，改从排列展开式的系数理解。
% 补齐递推初值、a=b情形、分块公式条件及叉积退化情形。
% 2026-10-05补证：核对北大《高等代数》第六版第二章，印刷页60--62（PDF页65--67）。
% 展开排列分组中两组内部求和的乘积，并补证叉积的零值判别。
\RequirePackage{iftex}
\RequireLuaTeX
\documentclass[UTF8, fontset=fandol, a4paper, 12pt]{ctexart}

\usepackage[margin=2.5cm]{geometry}
\usepackage{amsmath, amssymb, amsthm}
\usepackage{enumitem}
\usepackage{fancyhdr}
\usepackage[hidelinks]{hyperref}

\theoremstyle{definition}
\newtheorem{theorem}{定理}[section]
\newtheorem{proposition}[theorem]{命题}
\newtheorem{definition}[theorem]{定义}
\newtheorem{example}[theorem]{例题}
\renewcommand{\proofname}{证明}
\setlist[enumerate]{leftmargin=2em, itemsep=0.2em, topsep=0.3em}
\setlength{\headheight}{16pt}
\pagestyle{fancy}
\fancyhf{}
\fancyhead[L]{\small 高等代数（一）}
\fancyhead[R]{\small 行列式计算与拉普拉斯展开}
\fancyfoot[C]{\thepage}
\renewcommand{\headrulewidth}{0.3pt}
\raggedbottom

\title{行列式计算与拉普拉斯展开}
\author{Yuchen Zhang}
\date{课堂日期：2026年9月29日}
\makeatletter
\renewcommand{\maketitle}{%
  \begin{center}
    {\LARGE\bfseries \@title\par}\vspace{0.6em}
    \@author\par\vspace{0.2em}
    {\small \@date\par}
  \end{center}
}
\makeatother

\begin{document}
\maketitle

\section{从排列展开式理解代数余子式}

设 $A=(a_{ij})$ 为 $n$ 阶方阵（$n\ge2$），沿用上一讲的记号，
$M_{ij}$、$A_{ij}=(-1)^{i+j}M_{ij}$ 分别表示元素 $a_{ij}$ 的余子式、代数余子式。
固定第二列到第 $n$ 列，把第一列换成变量 $x_1,\ldots,x_n$，记
\[
 \Delta(x_1,\ldots,x_n)=
 \begin{vmatrix}x_1&a_{12}&\cdots&a_{1n}\\x_2&a_{22}&\cdots&a_{2n}\\
 \vdots&\vdots&&\vdots\\x_n&a_{n2}&\cdots&a_{nn}\end{vmatrix}.
\]
按列的次序写出排列定义，有
\[
 \Delta(x_1,\ldots,x_n)
 =\sum_{(i_1,\ldots,i_n)}(-1)^{\tau(i_1,\ldots,i_n)}
 x_{i_1}a_{i_2,2}\cdots a_{i_n,n},
\]
其中求和遍历 $1,\ldots,n$ 的所有排列，$\tau$ 表示逆序数。
每一项恰含一个 $x_k$。固定首个行号 $i_1=k$，它与其余行号形成 $k-1$ 个逆序。
将剩余行号按递增次序重新编号为 $1,\ldots,n-1$，不改变它们之间的逆序关系，
所以含 $x_k$ 的各项去掉这个因子后，其和为 $(-1)^{k-1}M_{k1}=A_{k1}$。
把含同一个 $x_k$ 的项合并，便得到
\[
 \Delta(x_1,\ldots,x_n)=\sum_{k=1}^n x_k A_{k1}.
\]
因此，代数余子式 $A_{k1}$ 就是把第一列看作变量时，$x_k$ 的系数。
它不依赖第一列中的元素，只由其余列决定。换成任意一列或一行，道理相同。

\subsection{不同行（列）的对应乘积之和}

上述观点也解释了上一讲的零和公式。若把第 $j$ 列换成第 $l$ 列（$l\ne j$），
新行列式有两列相同，故为零；第 $j$ 列的代数余子式却不变。因此
\[
 \sum_{i=1}^n a_{il}A_{ij}=0\quad(l\ne j).
\]
同理，按行有
\[
 \sum_{j=1}^n a_{kj}A_{ij}=0\quad(k\ne i).
\]
与按行（列）展开公式合在一起，即为
\[
 \sum_{r=1}^n a_{hr}A_{ir}
 =\begin{cases}\det A,&h=i,\\0,&h\ne i,\end{cases}
 \qquad
 \sum_{r=1}^n a_{rk}A_{rj}
 =\begin{cases}\det A,&k=j,\\0,&k\ne j.\end{cases}
\]
同一行（列）对应相乘求和得到原行列式，换成另一行（列）则得到零。

\section{行列式的降阶与递推计算}

\begin{example}
计算
\[
 D=\begin{vmatrix}
 5&3&-1&2&0\\1&7&2&5&2\\0&-2&3&1&0\\0&-4&-1&4&0\\0&2&3&5&0
 \end{vmatrix}.
\]
\end{example}
\begin{proof}[解]
先按第五列展开，再按所得四阶行列式的第一列展开，得
\[
 D=2(-1)^{2+5}
 \begin{vmatrix}5&3&-1&2\\0&-2&3&1\\0&-4&-1&4\\0&2&3&5\end{vmatrix}
 =-10\begin{vmatrix}-2&3&1\\-4&-1&4\\2&3&5\end{vmatrix}.
\]
对右侧三阶行列式作 $r_3\leftarrow r_3+r_1$，再从第三行提出 $6$，得
\[
 D=-60\begin{vmatrix}-2&3&1\\-4&-1&4\\0&1&1\end{vmatrix}
 =-60\bigl[(-2)(-1-4)-3(-4)+(-4)\bigr]
 =\boxed{-1080}.
\]
每次展开都选择非零元素最少的一行或一列；行倍加用于进一步制造零元素。
\end{proof}

\begin{example}
求主对角线为 $a+b$、上副对角线为 $a$、下副对角线为 $b$ 的三对角行列式
\[
 D_n=\begin{vmatrix}
 a+b&a&0&\cdots&0\\b&a+b&a&\ddots&\vdots\\
 0&b&a+b&\ddots&0\\\vdots&\ddots&\ddots&\ddots&a\\0&\cdots&0&b&a+b
 \end{vmatrix}\quad(n\ge1).
\]
\end{example}
\begin{proof}[解]
约定 $D_0=1$，则 $D_1=a+b$。按第一列展开，第二项的余子式再按第一行展开，得
\[
 D_n=(a+b)D_{n-1}-abD_{n-2}\quad(n\ge2).
\]
令 $E_n=D_n-bD_{n-1}$，则 $E_n=aE_{n-1}$，且 $E_1=a$，故 $E_n=a^n$。
于是 $D_n=a^n+bD_{n-1}$，逐次代入得
\[
 \boxed{D_n=\sum_{k=0}^n a^{n-k}b^k
 =\begin{cases}\dfrac{a^{n+1}-b^{n+1}}{a-b},&a\ne b,\\[6pt]
 (n+1)a^n,&a=b.\end{cases}}
\]
有限和的形式不需要除以 $a-b$，因而也适用于 $a=b$ 或参数为零的情形。
\end{proof}

\section{多行（列）的拉普拉斯展开}

\subsection{子式、余子式与展开公式}

设 $A$ 为 $n$ 阶方阵，$1\le r<n$。选取递增的行指标集
$I=\{i_1<\cdots<i_r\}$ 和列指标集 $J=\{j_1<\cdots<j_r\}$。
记 $A[I,J]$ 为这些行、列交叉位置构成的子矩阵，$I^c,J^c$ 为相应的补集。
所有子矩阵中的行、列均保持原来的顺序。

\begin{definition}
$r$ 阶行列式 $\det A[I,J]$ 称为一个 $r$ 阶\textbf{子式}。
删去所选行、列后所得的行列式，称为它的\textbf{余子式}，即 $\det A[I^c,J^c]$。
其\textbf{代数余子式}为
\[
 (-1)^{\sum I+\sum J}\det A[I^c,J^c],
 \qquad \sum I=i_1+\cdots+i_r,\quad\sum J=j_1+\cdots+j_r.
\]
\end{definition}

\begin{theorem}[拉普拉斯展开]
固定 $r$ 列 $J$，对所有 $r$ 行的选择 $I$ 求和，有
\[
 \boxed{\det A=\sum_{|I|=r}(-1)^{\sum I+\sum J}
       \det A[I,J]\det A[I^c,J^c].}
\]
固定 $r$ 行 $I$ 时，则对所有 $r$ 列的选择 $J$ 求和：
\[
 \det A=\sum_{|J|=r}(-1)^{\sum I+\sum J}
       \det A[I,J]\det A[I^c,J^c].
\]
\end{theorem}
每次固定一组行或列，只遍历另一组，共有 $\binom nr$ 项。
当 $r=1$ 时，这正是上一讲的按一行或一列展开公式。

\subsection{排列分组与展开式的符号}

\begin{proof}
固定 $J=\{1,\ldots,r\}$。排列展开式的每一项在前 $r$ 列取出的元素，
必定来自 $r$ 个不同的行。按这组行 $I=\{k_1<\cdots<k_r\}$ 把所有项分组。
再将其余行指标递增排列，记为 $I^c=\{l_1<\cdots<l_{n-r}\}$。

这一组中的每个完整排列都唯一写成
\[
 k_{\pi(1)},\ldots,k_{\pi(r)},
 l_{\rho(1)},\ldots,l_{\rho(n-r)},
 \qquad \pi\in S_r,\quad\rho\in S_{n-r}.
\]
其中 $\pi,\rho$ 分别决定两组内部的次序，可以彼此独立地选取。
完整排列的逆序分为两组内部的逆序和组间的逆序。
因为 $k_s,l_t$ 各按递增次序编号，两组内部的逆序数分别为 $\tau(\pi)$、$\tau(\rho)$。

所选行号总在前面，所以组间逆序只取决于 $I$，不取决于两组内部怎样排列。
对每个 $k_s$，比它小的行号共 $k_s-1$ 个，其中 $s-1$ 个在所选组内，
剩余的 $k_s-s$ 个均在后一组，因而各产生一个逆序。组间逆序数为
\[
 h_I=\sum_{s=1}^r(k_s-s)=\sum I-\frac{r(r+1)}2.
\]
因此完整排列的符号为 $(-1)^{h_I+\tau(\pi)+\tau(\rho)}$。
对所有 $\pi,\rho$ 求和，由分配律，这一组的和可以写成两个和的乘积：
\begin{align*}
 &(-1)^{h_I}
 \left(\sum_{\pi\in S_r}(-1)^{\tau(\pi)}
             \prod_{s=1}^r a_{k_{\pi(s)},s}\right)\\
 &\qquad\cdot
 \left(\sum_{\rho\in S_{n-r}}(-1)^{\tau(\rho)}
             \prod_{t=1}^{n-r}a_{l_{\rho(t)},r+t}\right)\\
 &=(-1)^{h_I}
   \det A[I,\{1,\ldots,r\}]\det A[I^c,\{r+1,\ldots,n\}].
\end{align*}
两个括号分别按排列定义计算子式、余子式。
每一对内部排列恰对应原展开式中的一项，故这一组共有 $r!(n-r)!$ 项。
不同的 $I$ 不会收进同一个排列，而每个排列都属于某个 $I$，
所以对所有 $I$ 求和，恰好得到原行列式的全部项。

对一般列集 $J=\{j_1<\cdots<j_r\}$，将这些列依次移到前面，且两组内部次序不变，
需要 $h_J=\sum_{s=1}^r(j_s-s)=\sum J-r(r+1)/2$ 次相邻对换。
设移列后所得矩阵为 $B$，则
\[
 \det B=(-1)^{h_J}\det A,\qquad \det A=(-1)^{h_J}\det B.
\]
在 $B$ 中应用前 $r$ 列的公式，两个小矩阵分别就是原矩阵的 $A[I,J]$、$A[I^c,J^c]$。
还原为 $\det A$ 时，每组再乘以 $(-1)^{h_J}$，所以总指数为
\[
 h_I+h_J=\sum I+\sum J-r(r+1).
\]
由于 $r(r+1)$ 是偶数，总符号为 $(-1)^{\sum I+\sum J}$，得到固定任意 $r$ 列的公式。
对转置矩阵应用该公式，就得到固定任意 $r$ 行的公式。
\end{proof}

\section{拉普拉斯展开的应用}

\subsection{分块三角行列式}

\begin{proposition}
设 $A$、$D$ 分别是 $p$ 阶、$q$ 阶方阵，$B$ 为 $p\times q$ 矩阵，则
\[
 \det\begin{pmatrix}A&B\\0&D\end{pmatrix}=\det A\det D.
\]
同理，$\det\begin{pmatrix}A&0\\C&D\end{pmatrix}=\det A\det D$。
\end{proposition}
\begin{proof}
按前 $p$ 列展开。这些列在后 $q$ 行全为零，因此只有选取前 $p$ 行时子式才可能非零。
该项的符号为 $(-1)^{2(1+\cdots+p)}=1$，子式为 $\det A$，余子式为 $\det D$。
下分块三角的情形由转置得到。
\end{proof}

这里要求对角块是方阵，但不要求它们可逆。一般的分块矩阵不能直接套用下面的二阶形式：
\[
 \det\begin{pmatrix}A&B\\C&D\end{pmatrix}
 =\det A\det D-\det B\det C.
\]
例如四个块均为 $2\times2$ 矩阵，取 $A=D=0$、$B=C=I_2$。
整体矩阵对应两次对换，行列式为 $1$，而右端为 $-1$。

\subsection{一个四阶行列式}

\begin{example}
从矩阵 $\begin{pmatrix}a_1&b_1&c_1\\a_2&b_2&c_2\end{pmatrix}$ 中分别删去第 $1,2,3$ 列，
所得二阶行列式依次记为
\[
 A=\begin{vmatrix}b_1&c_1\\b_2&c_2\end{vmatrix},\quad
 B=\begin{vmatrix}a_1&c_1\\a_2&c_2\end{vmatrix},\quad
 C=\begin{vmatrix}a_1&b_1\\a_2&b_2\end{vmatrix}.
\]
证明
\[
 H=\begin{vmatrix}a_1&b_1&c_1&0\\a_2&b_2&c_2&0\\
 0&a_1&b_1&c_1\\0&a_2&b_2&c_2\end{vmatrix}=AC-B^2.
\]
\end{example}
\begin{proof}
按前两行展开。含第四列的二阶子式均为零；选取第 $2,3$ 列时，余子式有一列为零。
故六种列选择中只有 $\{1,2\}$、$\{1,3\}$ 两项可能非零，分别为
\[
 (-1)^{1+2+1+2}CA=CA,\qquad
 (-1)^{1+2+1+3}BB=-B^2.
\]
相加即得 $H=AC-B^2$。此处 $A,B,C$ 均是二阶行列式的值，即三个数。
\end{proof}

\section{代数余子式与三维叉积}

设 $\boldsymbol\alpha=(a_{11},a_{21},a_{31})^{\mathsf T}$、
$\boldsymbol\beta=(a_{12},a_{22},a_{32})^{\mathsf T}$ 为 $\mathbb R^3$ 中的向量。
把它们作为前两列，第三列记为 $\boldsymbol x=(x,y,z)^{\mathsf T}$。按第三列展开：
\begin{align*}
 \det(\boldsymbol\alpha,\boldsymbol\beta,\boldsymbol x)
 &=\begin{vmatrix}a_{11}&a_{12}&x\\a_{21}&a_{22}&y\\a_{31}&a_{32}&z\end{vmatrix}\\
 &=x\begin{vmatrix}a_{21}&a_{22}\\a_{31}&a_{32}\end{vmatrix}
 -y\begin{vmatrix}a_{11}&a_{12}\\a_{31}&a_{32}\end{vmatrix}
 +z\begin{vmatrix}a_{11}&a_{12}\\a_{21}&a_{22}\end{vmatrix}.
\end{align*}
因此，若令
\[
 \boldsymbol\gamma=
 \begin{pmatrix}
 a_{21}a_{32}-a_{31}a_{22}\\
 a_{31}a_{12}-a_{11}a_{32}\\
 a_{11}a_{22}-a_{21}a_{12}
 \end{pmatrix},
\]
就有
\[
 \det(\boldsymbol\alpha,\boldsymbol\beta,\boldsymbol x)
 =\boldsymbol\gamma\cdot\boldsymbol x.
\]
取 $\boldsymbol x=\boldsymbol\alpha$ 或 $\boldsymbol\beta$，左侧都有两列相同，因而
\[
 \boldsymbol\gamma\cdot\boldsymbol\alpha=0,
 \qquad \boldsymbol\gamma\cdot\boldsymbol\beta=0.
\]
这解释了为什么由这三个代数余子式组成的向量，同时与原来的两个向量正交。

\begin{definition}
上述向量 $\boldsymbol\gamma$ 称为 $\boldsymbol\alpha$ 与 $\boldsymbol\beta$ 的\textbf{叉积}，记为
$\boldsymbol\alpha\times\boldsymbol\beta$。
\end{definition}

例如
\[
 (1,0,0)^{\mathsf T}\times(0,1,0)^{\mathsf T}=(0,0,1)^{\mathsf T}.
\]
这里称 $\boldsymbol\alpha,\boldsymbol\beta$ \textbf{线性相关}，是指存在不全为零的实数 $s,t$，
使 $s\boldsymbol\alpha+t\boldsymbol\beta=\boldsymbol0$；若只有 $s=t=0$ 才能使等式成立，则称为\textbf{线性无关}。

\begin{proposition}
$\boldsymbol\alpha\times\boldsymbol\beta=\boldsymbol0$ 当且仅当 $\boldsymbol\alpha,\boldsymbol\beta$ 线性相关。
\end{proposition}
\begin{proof}
若 $\boldsymbol\alpha=\boldsymbol0$，则叉积为零，且取 $(s,t)=(1,0)$ 即得一组不全为零的关系式系数。
以下设 $\boldsymbol\alpha\ne\boldsymbol0$。

若两向量线性相关，取一组不全为零的 $s,t$ 使关系式成立。
必有 $t\ne0$，否则 $s\boldsymbol\alpha=\boldsymbol0$ 会迫使 $s=0$。
于是 $\boldsymbol\beta=-(s/t)\boldsymbol\alpha$。
代入各二阶子式可知它们都为零，故叉积为零。

反过来，若叉积为零，由于它的三个坐标分别是三个二阶子式或其相反数，这些二阶子式均为零。
取 $\boldsymbol\alpha$ 的一个非零分量 $a_{p1}$，则对每个 $q=1,2,3$，有
\[
 a_{p1}a_{q2}-a_{q1}a_{p2}=0,
 \qquad a_{q2}=\frac{a_{p2}}{a_{p1}}a_{q1}.
\]
因此 $\boldsymbol\beta=(a_{p2}/a_{p1})\boldsymbol\alpha$，
取 $(s,t)=(-a_{p2}/a_{p1},1)$ 即得一组不全为零的关系式系数，所以两向量线性相关。
\end{proof}
线性无关的两向量具有非零叉积，它给出所张平面的一个法向量。
线性相关时叉积为零向量，不能用来指定法向方向。
交换两向量，相当于交换行列式的前两列，故
\[
 (\boldsymbol\beta\times\boldsymbol\alpha)\cdot\boldsymbol x
 =-(\boldsymbol\alpha\times\boldsymbol\beta)\cdot\boldsymbol x.
\]
分别取 $\boldsymbol x=\boldsymbol e_1,\boldsymbol e_2,\boldsymbol e_3$，就得到三个坐标均变号，
即 $\boldsymbol\beta\times\boldsymbol\alpha=-\boldsymbol\alpha\times\boldsymbol\beta$。
所以叉积的三个坐标并非另行记忆的一组公式，它们就是第三列的代数余子式，符号依次为 $+,-,+$。

\end{document}
